% ===================================================================
%  QUADRO N.º 2 — O Corpo Humano. — O Recreio. — Os Jogos.
%
%  Ceci n'est PAS une transcription : c'est une traduction, faite en
%  2026, en portugais d'aujourd'hui. Voir texto/pt/10-quadro-01.tex
%  pour la consigne, qui vaut pour les seize fichiers.
% ===================================================================

%%K t02-apar-1 apar
\VUsaut{4.0mm}
\VUtitre{15.0pt}{60}{QUADRO N.º 2}
\VUsaut{2.0mm}
\VUtitre{11.4pt}{0}{O Corpo Humano. --- O Recreio.}
\VUtitre{11.4pt}{0}{Os Jogos.}

%%K t02-tit-0 sub
\VUtitre{13.2pt}{0}{O Corpo Humano.}
\VUsaut{2.4mm}
\VUcentre{10.2pt}{0}{\textit{(Lição simples de ciências naturais.)}}
\VUsaut{3.0mm}

%%K t02-01-1 p
1. --- As principais partes do corpo humano são: a \VUgras{cabeça}, o
\VUgras{tronco} e os \VUgras{membros}.

%%K t02-tit-1 sub
\VUpk{10.2pt}{40}{A Cabeça.}
\VUsaut{3.0mm}

%%K t02-02-1 p
2. --- A cabeça consta de duas partes: o \VUgras{crânio}, que contém o
\VUgras{cérebro} e que o \VUgras{cabelo} cobre, e o \VUgras{rosto}.

%%K t02-03-1 p
3. --- No nosso desenho não vemos nem o crânio nem o cérebro; mas vemos
com toda a clareza as partes importantes do rosto.

%%K t02-04-1 p
4. --- Eis primeiro a \VUgras{testa}, que anuncia uma inteligência
poderosa, um pensamento sempre desperto. As \VUgras{têmporas} estão muito
desafogadas. O \VUgras{olho} é vivo e bem aberto. Uma \VUgras{sobrancelha}
espessa e umas \VUgras{pestanas} longas protegem-no.

%%K t02-05-1 p
5. --- Eis ainda o \VUgras{nariz} e as \VUgras{narinas}. O nariz serve-nos
para cheirar e para respirar. De cada lado do nariz estão as
\VUgras{faces}, e mais além as \VUgras{orelhas}. A parte inferior do rosto é
formada pela \VUgras{boca} e pelo \VUgras{queixo}.

%%K t02-06-1 p
6. --- Não as vemos muito bem, porque o \VUgras{bigode} e a \VUgras{barba}
no-las ocultam. Mas sabemos que a boca tem os dois \VUgras{lábios}, a
\VUgras{mandíbula superior} e a \VUgras{mandíbula inferior} com os seus
\VUgras{dentes}, e a \VUgras{língua}. Com a boca falamos e comemos. Com a
língua saboreamos os alimentos.

%%K t02-07-1 p
7. --- O olho, a orelha, o nariz e a boca são, juntamente com os dedos,
os órgãos dos cinco sentidos: a \VUgras{vista}, a \VUgras{audição}, o
\VUgras{olfato}, o \VUgras{paladar} e o \VUgras{tato}.

%%K t02-08-1 p
8. --- A cabeça é, pois, uma das partes mais interessantes do corpo
humano.

%%K t02-tit-2 sub
\VUsaut{4.4mm}
\VUpk{10.2pt}{40}{O Tronco.}
\VUsaut{3.0mm}

%%K t02-09-1 p
9. --- O \VUgras{pescoço} une a cabeça ao tronco. Compreende a
\VUgras{garganta} e a \VUgras{nuca}.

%%K t02-10-1 p
10. --- O tronco contém também muitos órgãos essenciais. No \VUgras{peito}
estão os \VUgras{pulmões}, com os quais respiramos, e o \VUgras{coração},
que é o centro da vida. Quando o coração deixa de bater, o ser vivo
morre.

%%K t02-11-1 p
11. --- As \VUgras{costelas} sustentam o peito.

%%K t02-12-1 p
12. --- As demais partes do tronco são o \VUgras{flanco}, as \VUgras{costas},
os \VUgras{ombros} e o \VUgras{abdome} (ventre). Deitamo-nos de lado ou de
costas para dormir, e levamos as cargas ao ombro.

%%K t02-13-1 p
13. --- No abdome estão o \VUgras{estômago} e os \VUgras{intestinos}. Servem
para digerir os alimentos.

%%K t02-tit-3 sub
\VUsaut{4.4mm}
\VUpk{10.2pt}{40}{Os Membros.}
\VUsaut{3.0mm}

%%K t02-14-1 p
14. --- Temos quatro \VUgras{membros}: dois \VUgras{braços} e duas
\VUgras{pernas}. Com os braços pegamos, seguramos, levantamos e apanhamos
os objetos; com as pernas mantemo-nos de pé, andamos e corremos.

%%K t02-15-1 p
15. --- O \VUgras{ombro} une o braço ao corpo. Um \VUgras{músculo} muito
robusto, o bíceps, move o braço, que o \VUgras{cotovelo} une ao
\VUgras{antebraço}. O \VUgras{pulso} forma a articulação da \VUgras{mão}, que
termina nos \VUgras{dedos} e nas \VUgras{unhas}. Na mão distinguimos uma
\VUgras{veia} (e uma artéria) por onde corre o \VUgras{sangue}.

%%K t02-16-1 p
16. --- As partes da perna são: a \VUgras{coxa}, o \VUgras{joelho} e a
\VUgras{curva do joelho}, a \VUgras{barriga da perna}, cuja \VUgras{carne} a
\VUgras{pele} cobre, o \VUgras{tornozelo} e o \VUgras{pé}. Os \VUgras{ossos} da
perna são muito grossos e muito fortes. Na extremidade do pé estão os
\VUgras{dedos do pé}. As demais partes são a \VUgras{planta} e o
\VUgras{calcanhar}.

%%K t02-17-1 p
17. --- Tal é a descrição quase completa do corpo humano.

%%K t02-17-2 p
\textit{Nota.} --- \textit{Com o auxílio da expressão «dói-me... (a
cabeça, etc.)», o professor poderá rever todo este vocabulário do ponto
de vista do bom ou do mau} \VUgras{estado do corpo}.

%%K t02-tit-4 sub
\VUsaut{5.0mm}
\VUpk{10.2pt}{40}{Ginástica.}
\VUsaut{2.4mm}
\VUcentre{10.2pt}{0}{\textit{(Contado por um dos alunos.)}}
\VUsaut{3.0mm}

%%K t02-18-1 p
18. --- Para sermos flexíveis, ágeis e sãos, temos de exercitar as pernas
e os braços. Por isso brincamos e fazemos \VUgras{ginástica}.

%%K t02-19-1 p
19. --- Durante a semana, estes dois exercícios fazem-se no pátio do
liceu (ver o quadro n.º 1). Mas às quintas-feiras e aos domingos vamos a
um grande prado, onde temos espaço para correr e divertir-nos.

%%K t02-20-1 p
20. --- Os nossos mestres e os nossos pais acompanham-nos lá às vezes;
mas é o \VUgras{professor de ginástica} \textsuperscript{(12)} quem dirige os
exercícios.

%%K t02-21-1 p
21. --- No prado, à esquerda da entrada, estão os \VUgras{aparelhos de
ginástica} \textsuperscript{(1)}: trepamos pela \VUgras{escada}
\textsuperscript{(2)}, ficamos de pernas para o ar nas \VUgras{argolas}
\textsuperscript{(3)} e no \VUgras{trapézio} \textsuperscript{(4)}, içamo-nos com as
mãos até ao alto da \VUgras{escada de corda} \textsuperscript{(5)} ou da
\VUgras{corda de nós} \textsuperscript{(6)}.

%%K t02-22-1 p
22. --- Entretanto, os nossos companheiros saltam por cima do \VUgras{cavalo
de madeira} \textsuperscript{(7)} ou das \VUgras{barras paralelas}
\textsuperscript{(11)}. Outros \VUgras{boxeiam} \textsuperscript{(8)} ou fazem
\VUgras{esgrima} \textsuperscript{(9)} com máscara e \VUgras{florete}. Outros,
enfim, levantam \VUgras{pesos} \textsuperscript{(10)}.

%%K t02-tit-5 sub
\VUsaut{4.6mm}
\VUpk{10.2pt}{40}{Os Jogos.}
\VUsaut{3.0mm}

%%K t02-23-1 p
23. --- À volta dos ginastas, meninos e meninas jogam diferentes
\VUgras{jogos}. As meninas divertem-se com o seu \VUgras{arco}
\textsuperscript{(14)}; saltam à \VUgras{corda} \textsuperscript{(13)}, ou convidam os
seus irmãos e os seus primos para uma partida de \VUgras{croqué}
\textsuperscript{(15)}. Adoram afastar a \VUgras{bola} do adversário com o seu
\VUgras{maço de madeira}, justamente quando este julgava passar sem
trabalho pelo arco!

%%K t02-24-1 p
24. --- Os meninos preferem os jogos mais rudes. Jogam de bom grado ao
\VUgras{boliche} \textsuperscript{(16)}, cujos paus derrubam com manha com a
\VUgras{bola} \textsuperscript{(17)}. Também não desdenham o \VUgras{pião}
\textsuperscript{(18)} e o \VUgras{bilboquê} \textsuperscript{(19)}, nem sequer o
\VUgras{jogo da rã} \textsuperscript{(20)}. Mas os seus jogos preferidos são as
\VUgras{bolinhas de gude} \textsuperscript{(21)} e a \VUgras{pelota à mão}
\textsuperscript{(22)}, porque a parede da \VUgras{estufa} \textsuperscript{(26)} serve
muito bem para fazer ressaltar a bola leve.

%%K t02-25-1 p
25. --- O pequeno João, empoleirado nas suas \VUgras{andas}
\textsuperscript{(24)}, é muito destro e sabe guardar muito bem o equilíbrio.
Perto dele, alguns companheiros brincam às \VUgras{escondidas}
\textsuperscript{(25)} nessa estufa e atrás da \VUgras{sebe}. Outros preferem a
\VUgras{cabra-cega com palmada} \textsuperscript{(28)} ou a \VUgras{peteca}
\textsuperscript{(29)}, que esvoaça de uma \VUgras{raquete} para a outra.

%%K t02-noto-1 noto
\VUnotes{143.2mm}{%
(*) Os jogos infantis têm muitas vezes nomes puramente nacionais.
Nomeámo-los em ido o melhor que pudemos, ora inspirando-nos na ação
mesma que os caracteriza, ora adotando o nome nacional de um ou outro
país quando não resulta demasiado inexato. Por exemplo: o \VUgras{jogo do
tonel} (alemão e francês) é o \VUgras{jogo da rã} \textsuperscript{(20)}
(inglês), no qual os jogadores tratam de deitar os seus discos redondos
pela boca da rã de metal que está de boca aberta sobre a caixa (o tonel)
furada. O \VUgras{salto ao ombro} \textsuperscript{(30)} distingue-se do
\VUgras{salto da rã} só nisto: para saltar por cima da cabeça, os
jogadores apoiam as mãos nos ombros do companheiro, ao passo que no
«\VUgras{salto da rã}» as apoiam só nas suas costas. --- Ou então, alguns
dos que jogam inclinam-se enquanto os outros saltam sobre as suas costas
apoiando as mãos nos ombros; os primeiros devem aguentar o golpe sem
dobrar-se, os segundos devem saltar sem perder o equilíbrio (ver o
próprio quadro).%
}

%%K t02-26-1 p
26. --- No meio do prado há duas árvores. Ao pé de uma delas, sete ou
oito rapazes organizaram uma partida de \VUgras{salto ao ombro}
\textsuperscript{(30)} (*). Este jogo não se conhece em todos os países; mas
toda a gente sabe o que é o \VUgras{salto da rã}, que os meninos não estão
a jogar agora.

%%K t02-tit-6 sub
\VUsaut{5.0mm}
\VUpk{10.2pt}{40}{Os Jogos ao ar livre.}
\VUsaut{3.4mm}

%%K t02-27-1 p
27. --- A \VUgras{cabra-cega} \textsuperscript{(31)} é também muito divertida;
meninas e meninos põem à vez a \VUgras{venda} nos olhos e apanham os
desajeitados que se deixam agarrar.

%%K t02-28-1 p
28. --- No meio do prado, eis um jogo muito francês, ainda que um tanto
brusco: é a \VUgras{ursa} \textsuperscript{(32)}, muito mais ruidosa do que o
jogo tão sossegado do \VUgras{papagaio de papel} \textsuperscript{(33)} ou do
que o jogo inglês do \VUgras{ténis} \textsuperscript{(34)}.

%%K t02-29-1 p
29. --- Este último desporto faz as delícias dos alunos mais velhos. Os
nossos próprios professores praticam-no de bom grado. Exige muita
destreza e agilidade, e eu gosto imenso dele. Deixo o \VUgras{baloiço}
\textsuperscript{(35)} para as meninas.

%%K t02-30-1 p
30. --- Também não gosto de outro jogo inglês que pode chegar a ser
perigoso.

%%K t02-30-1b p
Refiro-me ao \VUgras{futebol} \textsuperscript{(36)}, que faz a felicidade do meu
irmão mais velho. Prefiro o nosso velho \VUgras{jogo das barras}
\textsuperscript{(38)}, igualmente são e muito menos brutal.

%%K t02-tit-7 sub
\VUsaut{4.6mm}
\VUpk{10.2pt}{40}{Ciclismo e jogos de bola.}
\VUsaut{3.0mm}

%%K t02-31-1 p
31. --- Também dou com gosto uma volta ao prado de \VUgras{bicicleta}
\textsuperscript{(37)}.

%%K t02-32-1 p
32. --- No ano que vem aprenderei a \VUgras{montar a cavalo}
\textsuperscript{(39)}. Que belo é um cavalo que anda ou trota com garbo, e
que salta os obstáculos a galope!

%%K t02-33-1 p
33. --- No verão aprendemos a \VUgras{nadar} \textsuperscript{(40)}. Mergulhamos e
banhamo-nos com deleite na água clara e fresca do rio. Às vezes, no fim,
soltamos um \VUgras{balão de papel} \textsuperscript{(41)} que sobe
majestosamente pelo ar assim que se enche de ar quente.

%%K t02-34-1 p
34. --- Há ainda muitos outros jogos, mas nunca acabaria de nomeá-los, e
por este breve resumo vê-se que as quintas-feiras no nosso belo prado não
são nada aborrecidas.

%%K t02-34-2 p
N. B. --- \textit{O professor completará facilmente este capítulo
descrevendo com mais pormenor cada um dos jogos importantes.}
\VUsaut{6.0mm}
\VUfilet{22mm}
